\begin{ouverture}[synthese]{9}{Synthétiser et identifier}{Constitution et transformations}{Observer · Modéliser · Expérimenter · Expliquer}
\objectif{Expliquer l’intérêt d’un chauffage à reflux.}
\objectif{Légender les montages nécessaires à une synthèse.}
\objectif{Comparer une espèce synthétisée à une espèce naturelle.}
\prerequis{Proportionnalité, conversions et lecture de graphiques.}
\end{ouverture}
\essentiel
\essSection{1}{Une espèce et son origine}
Une synthèse fabrique une espèce à partir d'autres espèces par transformation chimique. Une espèce présente dans la nature peut être préparée au laboratoire. Si les molécules sont identiques, leurs propriétés intrinsèques sont les mêmes, quelle que soit l'origine ; les échantillons peuvent néanmoins différer par leurs impuretés.
\begin{attention}[Naturel ne signifie pas sans danger]L'origine d'un échantillon ne suffit pas à juger sa dangerosité, sa pureté ou son impact environnemental. Lire les données et les pictogrammes de sécurité de chaque produit.\end{attention}
\essSection{2}{Chauffer sans perdre le milieu}
Le chauffage à reflux chauffe un mélange tout en condensant les vapeurs dans un réfrigérant vertical. Le liquide condensé retombe dans le ballon : on limite les pertes d'espèces volatiles. Le réfrigérant reste ouvert vers l'atmosphère ; l'eau entre par le bas et sort par le haut.
\schema[corrige]{reflux}
Le montage est maintenu par un support. Des grains anti-bouillonnement peuvent régulariser l'ébullition si le protocole le prévoit ; ils se mettent avant le chauffage. Le reflux ne sépare pas, à lui seul, le produit des autres espèces.
\essSection{3}{Isoler et contrôler}
Après réaction, le protocole peut prévoir une séparation, un lavage ou une filtration, puis une identification. Une chromatographie compare produit synthétisé et références sur la même plaque, dans le même éluant.
\schema[corrige]{ccm-synthese}
\begin{savoirfaire}[Argumenter]Dire quelles taches se correspondent ; rechercher des taches supplémentaires ; conclure prudemment sur l'identité et la pureté dans les limites de la méthode. Une seule tache ne prouve pas l'absence de toute impureté non détectée.\end{savoirfaire}
\essSection{4}{Expérimenter et modéliser}
\begin{experience}[Concevoir un compte rendu de synthèse]À partir d'un protocole fourni par l'enseignant, préparer le schéma légendé du reflux et la fiche de sécurité, identifier les étapes de transformation, d'isolement et de contrôle. Comparer ensuite par CCM le produit obtenu et une référence. Ne réaliser aucune synthèse domestique à partir du seul schéma du livre.\end{experience}
\exercices
\rubrique{Prendre confiance}
\begin{blocLivre}{ex-9-1}\begin{exercice}[1]{L’origine d’une molécule}
Une molécule de vanilline synthétisée a la même structure que celle extraite de la vanille.
\begin{questions}
\item Constituent-elles la même espèce chimique ?
\item Les deux échantillons sont-ils nécessairement de même composition ?
\end{questions}
\end{exercice}\end{blocLivre}
\begin{blocLivre}{ex-9-2}\begin{exercice}[1]{Reconnaître le reflux}
Observer le montage fourni.
\schema{reflux}
\begin{questions}
\item Nommer le ballon, le réfrigérant et le chauffage.
\item Indiquer le sens de circulation de l’eau.
\end{questions}
\end{exercice}\end{blocLivre}
\begin{blocLivre}{ex-9-3}\begin{exercice}[1]{Le rôle du réfrigérant}
Un mélange contient un solvant volatil et doit être chauffé.
\begin{questions}
\item Pourquoi utiliser un reflux ?
\item Le réfrigérant empêche-t-il toute transformation chimique ?
\end{questions}
\end{exercice}\end{blocLivre}
\begin{blocLivre}{ex-9-4}\begin{exercice}[1]{Une cuve de CCM}
La ligne de dépôt est située à 1 cm du bas d’une plaque. L’éluant atteint 0,5 cm dans la cuve.
\begin{questions}
\item La position du dépôt convient-elle ?
\item Que risquerait-on si l’éluant atteignait 1,5 cm ?
\end{questions}
\end{exercice}\end{blocLivre}
\rubrique{S’entraîner}
\begin{blocLivre}{ex-9-5}\begin{exercice}[2]{Comparer deux produits}
Les pistes N et S représentent une référence naturelle et un échantillon synthétisé ; R est un réactif de départ.
\schema{ccm-synthese}
\begin{questions}
\item Quel indice associe S à l’espèce N ?
\item Quel indice révèle une impureté de S ?
\end{questions}
\end{exercice}\end{blocLivre}
\begin{blocLivre}{ex-9-6}\begin{exercice}[2]{Un montage à corriger}
Un élève ferme hermétiquement le haut du réfrigérant et branche l’arrivée d’eau en haut.
\begin{questions}
\item Corriger les deux choix.
\item Pourquoi le premier présente-t-il un risque ?
\end{questions}
\end{exercice}\end{blocLivre}
\begin{blocLivre}{ex-9-7}\begin{exercice}[2]{Étapes d’un protocole}
Un protocole indique : chauffer à reflux, laisser refroidir, filtrer, analyser le solide par CCM après dissolution adaptée.
\begin{questions}
\item Associer transformation, isolement et contrôle à ces étapes.
\item Pourquoi le refroidissement précède-t-il la manipulation suivante ?
\end{questions}
\end{exercice}\end{blocLivre}
\begin{blocLivre}{ex-9-8}\begin{exercice}[2]{Une seule tache}
Un produit présente une seule tache visible sur une CCM.
\begin{questions}
\item Quelle conclusion limitée peut-on formuler ?
\item Peut-on affirmer une pureté absolue ?
\end{questions}
\end{exercice}\end{blocLivre}
\rubrique{Raisonner et approfondir}
\begin{blocLivre}{ex-9-9}\begin{exercice}[3]{Un test d’identité incomplet}
Une espèce synthétisée et une référence ont la même masse volumique à la précision de la mesure.
\begin{questions}
\item Est-ce une preuve suffisante ?
\item Proposer deux contrôles complémentaires.
\end{questions}
\end{exercice}\end{blocLivre}
\begin{blocLivre}{ex-9-10}\begin{exercice}[3]{Légender utilement}
Un schéma de reflux ne porte que les noms « tube » et « pot ».
\schema{reflux}
\begin{questions}
\item Pourquoi les légendes sont-elles insuffisantes ?
\item Quelles informations de fonctionnement faut-il ajouter ?
\end{questions}
\end{exercice}\end{blocLivre}
\begin{blocLivre}{ex-9-11}\begin{exercice}[3]{Une comparaison équitable}
Deux CCM ont été réalisées dans des éluants différents. Une tache a la même hauteur sur les deux plaques.
\begin{questions}
\item Peut-on en déduire une identité ?
\item Proposer une comparaison plus fiable.
\end{questions}
\end{exercice}\end{blocLivre}
\begin{blocLivre}{ex-9-12}\begin{exercice}[3]{Limiter l’impact}
Deux protocoles produisent la même espèce en quantité comparable. A utilise 100 mL d’un solvant à risque ; B en utilise 20 mL d’un solvant moins dangereux selon les fiches fournies.
\begin{questions}
\item Quel choix paraît préférable sur ces critères ?
\item Pourquoi faut-il garder d’autres critères dans l’analyse ?
\end{questions}
\end{exercice}\end{blocLivre}

